\begin{table}[htbp]
\centering
\caption{Standardized Effect Sizes}
\label{tab:sde}
\begin{threeparttable}
\small
\begin{tabular}{lcccccc}
\toprule
Outcome & $\hat{\beta}$ & SE & SD(Y) & SDE & SE(SDE) & Classification \\
\midrule
\multicolumn{7}{l}{\textit{Panel A: Pooled}} \\
\addlinespace
Divorce rate & -0.035 & 0.160 & 0.771 & -0.046 & 0.207 & Small negative \\
Marriage rate & -0.369 & 0.765 & 1.647 & -0.224 & 0.465 & Large negative \\
\addlinespace
\multicolumn{7}{l}{\textit{Panel B: Heterogeneous (Early vs.\ Late Adopters)}} \\
\addlinespace
Divorce rate (early adopters) & 0.400 & NA & 1.345 & 0.297 & NA & Large positive \\
Divorce rate (late adopters) & 0.066 & 0.480 & 0.717 & 0.091 & 0.669 & Moderate positive \\
\bottomrule
\end{tabular}
\begin{tablenotes}[flushleft]
\small
\item \textit{Notes:} \textbf{Country:} United States. \textbf{Research question:} Do state policies that reduce marriage license fees for couples completing premarital counseling lower divorce rates? \textbf{Policy mechanism:} Ten states adopted policies offering \$20--60 marriage license fee reductions or waiting period waivers to couples completing 4--8 hours of premarital education courses, creating a financial incentive to acquire relationship skills before marriage. \textbf{Outcome definition:} State-level divorce rate per 1,000 population from CDC National Vital Statistics System. \textbf{Treatment:} Binary indicator for state adoption of a premarital education promotion policy. \textbf{Data:} CDC NVSS state divorce and marriage rates, 1990--2022, state-year level, approximately 1,500 state-year observations (excluding Georgia due to missing data). \textbf{Method:} Callaway-Sant'Anna (2021) doubly-robust estimator with never-treated controls; wild cluster bootstrap for inference with 10 treated clusters. \textbf{Sample:} 50 US states plus DC, excluding Georgia (missing divorce data 2004--2016); 9 treated states with staggered adoption 1998--2018, approximately 41 never-treated controls. SDE $= \hat{\beta} / \text{SD}(Y)$ where SD($Y$) is the pre-treatment standard deviation. Classification refers to magnitude, not statistical significance: Large ($|$SDE$| > 0.15$), Moderate ($0.05$--$0.15$), Small ($0.005$--$0.05$), Null ($< 0.005$).
\end{tablenotes}
\end{threeparttable}
\end{table}
